\section{Template}
\begin{frame}{Template}

\begin{itemize}
\item I \emph{template} invece possiamo pensarli come delle raccolte di comandi
che definiscono l'aspetto e i possibili comandi per un certo documento
\item Possono definire anche i pacchetti utilizzati
\item Non definiremo mai un template ma ce ne sono moltissimi disponibili online
    \subsection{Esempio di template per la tesi}
\item \textbf{Vediamo un esempio di template per la tesi!}
\end{itemize}

\begin{figure}
	\centering
	\includegraphics[scale=0.20]{res/images/template}
\end{figure}

\end{frame}
